 
\documentclass[../../../physics-standard_answers_all.tex]{subfiles}

\begin{document}

\setlength{\abovedisplayskip}{2mm}
\setlength{\belowdisplayskip}{1mm}

\begin{enumerate}

    \item R$_2$の電流を$i$（右向き）とすると，R$_3$の電流は$I-i$（右向き）である．
        キルヒホッフ則より，
        %
        \begin{align*}
            \left\{
                \begin{aligned}
                    & 2Ri + 2RI = 3V_0 \; , \\
                    & 2Ri = R(I-i) + V_0
                \end{aligned}
            \right.
            \qquad \text{∴} \; i = \frac{5V_0}{8R} \;,
            \quad I = \uwave{\frac{7V_0}{8R}} \;.
        \end{align*}

    \item R$_1$の電流を$i$（右向き），R$_5$の電流を$j$（上向き）とすると，
        R$_2$の電流は$i+j$（右向き），R$_3$の電流は$I-i$（右向き），
        R$_4$の電流は$I-i-j$（右向き）である．キルヒホッフ則より，
        %
        \begin{align*}
            & \left\{
                \begin{aligned}
                    & R(i+j) + Ri = 3R(I-i-j) + R(I-i) = V_0 \; , \\
                    & Ri = Rj + R(I-i)
                \end{aligned}
            \right. \\
            & \text{∴} \; i = \frac{4V_0}{9R} \;,
            \quad j = \frac{V_0}{9R} \;,
            \quad I = \uwave{\frac{7V_0}{9R}} \;.
        \end{align*}
    
\end{enumerate}

\end{document}
